\begin{frame}{Sincronizador de dois flip-flops}
  Para levar um sinal de um domínio de clock para outro, é comum usar dois flip-flops D em série.
  \begin{center}
  \resizebox{0.9\linewidth}{!}{%
  \begin{tikzpicture}[every node/.style={align=center}, node distance=1.1cm, >=Latex]
    \node (in) {Sinal assíncrono};
    \node[draw, minimum width=1.4cm, minimum height=1.2cm, right=of in] (ff1) {FF D\\1};
    \node[draw, minimum width=1.4cm, minimum height=1.2cm, right=of ff1] (ff2) {FF D\\2};
    \node[right=of ff2] (out) {Sinal sincronizado};
    \draw[->] (in)--(ff1); \draw[->] (ff1)--(ff2); \draw[->] (ff2)--(out);
    \draw[->, blue] (ff1.south) -- ++(0,-0.8) node[below]{CLK de destino};
    \draw[->, blue] (ff2.south) -- ++(0,-0.8);
  \end{tikzpicture}%
  }
  \end{center}
  \begin{alertblock}{Atenção}
    Esse circuito é indicado para sinais de controle de um bit. Para barramentos, usam-se protocolos, handshake ou FIFO.
  \end{alertblock}
\end{frame}